\documentclass[10pt]{article}
\usepackage{amsmath,amssymb}
\usepackage{longtable}
\usepackage[a4paper,left=2cm,right=2cm,top=1.5cm,bottom=2cm]{geometry}
\setlength{\emergencystretch}{2em}

\setlength{\parindent}{0pt}

\begin{document}

\title{Forward and Reverse Automatic Differentiation by Hand}
\author{}
\date{}
\maketitle
\subsection*{Learning goals}
After this exercise, you should be able to

\begin{itemize}
\item
explain what a \textbf{\textit{dual number}} is and why dual numbers are useful for forward-mode automatic differentiation,

\item
propagate values and derivatives through a computation using dual numbers,

\item
compute directional derivatives and Jacobian columns with forward mode,

\item
explain the idea of a \textbf{\textit{computational graph}} and an \textbf{\textit{adjoint}} in reverse-mode automatic differentiation,

\item
propagate adjoints backwards through a computation graph,

\item
compute gradients and Jacobian rows with reverse mode,

\item
explain the practical difference between forward and reverse AD.

\end{itemize}
The calculations in this sheet are designed to be done \textbf{\textit{by hand}}. No programming is required.

\clearpage
\section*{Part I - Forward-mode AD with dual numbers}
\subsection*{1. What are dual numbers?}
Automatic differentiation (AD) computes derivatives by applying the chain rule to the elementary operations of a calculation. For a composition, the chain rule gives $(f\circ q)'(x)=f'(q(x))q'(x)$.

A \textbf{\textit{dual number}} $\hat{x}$ has the form

\[
\hat{x} = x + \dot{x}\varepsilon
\]

where $x,\dot{x}\in\mathbb{R}$ and $\varepsilon$ is a formal symbol with the special property

\[
\varepsilon^2 = 0
\qquad
\varepsilon \neq 0
\]

The ordinary value $x$ is called the \textbf{\textit{primal value}}. The coefficient $\dot{x}$ is the \textbf{\textit{tangent}}: it tracks the derivative along a chosen input direction. The dot is notation for this derivative component, not necessarily a time derivative. Choosing the initial tangent is called setting an \textbf{\textit{input seed}}. \textbf{\textit{Forward mode}} carries primal values and tangents together from inputs to outputs.

If we then evaluate a differentiable function $f$ using its value and first derivative for each elementary operation together with $\varepsilon^2=0$, we obtain

\[
f(x+\dot{x}\varepsilon)
= f(x) + f'(x)\dot{x}\,\varepsilon
\]

For the special choice $\dot{x}=1$

\[
f(x+\varepsilon)=f(x)+f'(x)\varepsilon
\]

so the coefficient of $\varepsilon$ is directly $f'(x)$.

\subsection*{2. Arithmetic with dual numbers}
Let

\[
\hat{a}=a+\dot{a}\varepsilon,
\qquad
\hat{b}=b+\dot{b}\varepsilon.
\]

Then

\subsubsection*{Addition}
\[
\hat{a}+\hat{b}
=(a+b)+(\dot{a}+\dot{b})\varepsilon.
\]

\subsubsection*{Multiplication}
\[
\hat{a}\hat{b}
=(a+\dot{a}\varepsilon)(b+\dot{b}\varepsilon).
\]

Expanding the product gives

\[
ab+(a\dot{b}+\dot{a}b)\varepsilon
\]

which is exactly the product rule.

\subsubsection*{Some useful elementary functions}
Here $\log$ denotes the natural logarithm.

\[
\exp(a+\dot{a}\varepsilon)
= e^a+\dot{a}\,e^a\varepsilon,
\]

\[
\sin(a+\dot{a}\varepsilon)
= \sin(a)+\dot{a}\,\cos(a)\varepsilon,
\]

\[
\log(a+\dot{a}\varepsilon)
= \log(a)+\frac{\dot{a}}{a}\varepsilon,
\qquad a>0,
\]

and

\[
(a+\dot{a}\varepsilon)^2
= a^2+2a\,\dot{a}\varepsilon.
\]

\subsection*{3. Very simple 1D example}
Consider

\[
f(x)=x^2+3x
\]

and compute $f(2)$ and $f'(2)$ using dual numbers.

The dual $\hat{x}_0$ is (with tangent $\dot{x}_0=1$)

\[
\hat{x}_0=2+\varepsilon.
\]

Then

\[
\hat{x}_0^2=(2+\varepsilon)^2
=4+4\varepsilon,
\]

and

\[
3\hat{x}_0=6+3\varepsilon.
\]

Therefore

\[
f(\hat{x}_0)
=10+7\varepsilon.
\]

Hence

\[
f(2)=10
\]
\[
f'(2)=7.
\]

The important point is that the value and the derivative were propagated at the same time.

\subsection*{4. More than one input}
A scalar is a single number; a vector is an ordered collection of numbers. Vectors are written as column vectors, and $T$ denotes transposition.

Suppose now that

\[
g:\mathbb{R}^n\to\mathbb{R}^m.
\]

We can seed every input with its own tangent component:

\[
\hat{\mathbf{x}}
=\mathbf{x}+\dot{\mathbf{x}}\varepsilon.
\]

Forward AD then gives

\[
{
 g(\mathbf{x}+\dot{\mathbf{x}}\varepsilon)
 = g(\mathbf{x}) + J_g(\mathbf{x})\,\dot{\mathbf{x}}\,\varepsilon
}
\]

The \textbf{\textit{Jacobian matrix}} $J_g(\mathbf{x})$ collects all first partial derivatives: its entry in row $i$ and column $j$ is $\partial g_i/\partial x_j$. A \textbf{\textit{partial derivative}} measures the change with respect to one input while holding the others fixed. For $n$ inputs and $m$ outputs, the Jacobian has $m$ rows and $n$ columns.

A \textbf{\textit{sweep}} is one traversal of the calculation in the chosen direction. One forward-mode sweep computes one \textbf{\textit{Jacobian-vector product (JVP)}}

\[
J_g(\mathbf{x})\dot{\mathbf{x}}.
\]

A \textbf{\textit{standard basis vector}} has one entry equal to $1$ and all other entries equal to $0$. If we choose $\dot{\mathbf{x}}$ to be a standard basis vector, for example

\[
\mathbf{e}_1=(1,0,0)^T,
\]

then the tangent part is the first column of the Jacobian.

For a scalar-valued function, the \textbf{\textit{gradient}} $\nabla g$ is the column vector of all partial derivatives. The \textbf{\textit{directional derivative}} along an input direction $\mathbf v$ is

\[
D_{\mathbf v}g(\mathbf x)=\left.\frac{d}{dt}g(\mathbf x+t\mathbf v)\right|_{t=0}=\nabla g(\mathbf x)^T\mathbf v.
\]

Here $\mathbf v$ need not have unit length: scaling the seed scales the directional derivative.

\section*{Forward-mode exercises}
\subsection*{Exercise 1.1 - Scalar input, scalar output}
Consider

\[
f(x)=\log\left(xe^x+3\right).
\]

Evaluate the function and its derivative at

\[
x=1
\]

using dual-number propagation.

\subsection*{Exercise 1.2 - Vector input, scalar output}
Consider

\[
g(x,y,z)=xy+\sin z+y^2
\]

at the point

\[
(x,y,z)=(1,2,0).
\]

\begin{enumerate}
\item
First calculate the ordinary function value $g(1,2,0)$.

\item
Use forward-mode AD with the three seeds

\[
\mathbf{e}_x=(1,0,0),
   \qquad
   \mathbf{e}_y=(0,1,0),
   \qquad
   \mathbf{e}_z=(0,0,1)
\]

to determine the full gradient

\[
\nabla g(1,2,0).
\]

\item
Now use only one forward sweep with seed direction

\[
\dot{\mathbf{x}}=(1,-1,2)^T
\]

to compute the directional derivative

\[
D_{\dot{\mathbf{x}}}g
   =
   \nabla g^T\dot{\mathbf{x}}.
\]

\end{enumerate}
\subsection*{Exercise 1.3 - Vector input, vector output}
Consider the vector-valued function

\[
\mathbf{h}(x,y,z)
=
\begin{pmatrix}
 h_1(x,y,z)\\
 h_2(x,y,z)
\end{pmatrix}
=
\begin{pmatrix}
 xy+z\\
 x^2+\sin y-z^2
\end{pmatrix}.
\]

Evaluate it at

\[
(x,y,z)=(1,0,1).
\]

\begin{enumerate}
\item
Calculate $\mathbf h(1,0,1)$.

\item
Use the three input basis directions

\[
\mathbf e_x,
   \qquad
   \mathbf e_y,
   \qquad
   \mathbf e_z
\]

and dual numbers to compute the full Jacobian

\[
J_{\mathbf h}
   =
   \begin{pmatrix}
   \frac{\partial h_1}{\partial x} &
   \frac{\partial h_1}{\partial y} &
   \frac{\partial h_1}{\partial z}\\[4pt]
   \frac{\partial h_2}{\partial x} &
   \frac{\partial h_2}{\partial y} &
   \frac{\partial h_2}{\partial z}
   \end{pmatrix}.
\]

Remember: each forward sweep gives one column of the Jacobian.

\item
Use a single forward sweep with

\[
\mathbf v=(1,2,-1)^T
\]

to compute

\[
J_{\mathbf h}\mathbf v
\]

\end{enumerate}
\clearpage
\section*{Part II - Reverse-mode AD / Backward AD}
\subsection*{1. The basic idea}
Forward mode propagates derivatives together with the values from input to output.

A \textbf{\textit{computational graph}} represents a calculation as nodes for inputs and intermediate results, with directed edges showing which values each operation uses. An \textbf{\textit{intermediate variable}} stores the result of one of these operations. A \textbf{\textit{local derivative}} is the derivative of an operation's output with respect to one of its inputs.

\textbf{\textit{Reverse mode}} proceeds as follows:

\begin{enumerate}
\item
Perform an ordinary forward pass and store the intermediate values.

\item
Start from the output.

\item
Propagate sensitivities backwards through the computational graph.

\end{enumerate}
For an intermediate variable $v$, define its \textbf{\textit{adjoint}} (the sensitivity of the chosen scalar output to this variable) as

\[
{
\bar v = \frac{\partial L}{\partial v}
}
\]

where $L$ is the final scalar output whose derivative we want (the bar notation is common in reverse-mode AD).

Initialize all adjoints to zero. The \textbf{\textit{output seed}} specifies the starting adjoint at the output. For the derivative of the scalar output itself, set

\[
\bar L
=
\frac{\partial L}{\partial L}
=1.
\]

Then the chain rule is applied locally, one operation at a time, in reverse order.

\subsection*{2. Local backward rules}
Suppose an intermediate variable $c$ is computed from earlier variables.

\subsubsection*{Addition}
If

\[
c=a+b,
\]

then

\[
\bar a \mathrel{+}= \bar c,
\qquad
\bar b \mathrel{+}= \bar c.
\]

\textbf{\textit{Accumulation}}, written $\bar a\mathrel{+}=\bar c$, means replacing $\bar a$ by its current value plus $\bar c$: if one variable influences the output through several paths, all contributions to its adjoint must be added.

\subsubsection*{Multiplication}
If

\[
c=ab,
\]

then

\[
\bar a \mathrel{+}= \bar c\,b,
\qquad
\bar b \mathrel{+}= \bar c\,a.
\]

\subsubsection*{Some useful elementary functions}
\[
c = e^a \qquad \Rightarrow \qquad \bar a \mathrel{+}= \bar c\,e^a
\]
\[
c=\sin a \qquad \Rightarrow \qquad \bar a \mathrel{+}= \bar c\cos a.
\]
\[
c=\log a \qquad \Rightarrow \qquad \bar a \mathrel{+}= \bar c\frac{1}{a}.
\]
\[
c=a^2 \qquad \Rightarrow \qquad \bar a \mathrel{+}= \bar c\,2a.
\]

\subsection*{3. Very simple 1D reverse-mode example}
Consider

\[
f(x)=(x+1)^2
\]

at $x=2$.

Break the function into elementary operations:

\[
v_1=x+1,
\]

\[
v_2=v_1^2,
\]

\[
f=v_2.
\]

\subsubsection*{Forward pass}
At $x=2$,

\[
v_1=3,
\qquad
v_2=9.
\]

\subsubsection*{Backward pass}
Start with

\[
\bar v_2=1
\]
because $f=v_2$, so $\partial f/\partial v_2=1$.

Since

\[
v_2=v_1^2
\]

we get

\[
\bar v_1
=
\bar v_2\,2v_1
=1\cdot 6
=6
\]

and since

\[
v_1=x+1
\]

we obtain

\[
\bar x=\bar v_1=6
\]

Hence

\[
f'(2)=6.
\]

In one dimension this may look more complicated than ordinary differentiation. Its advantage becomes clear when a function has many inputs but only one scalar output.

\subsection*{4. Reverse mode for several inputs}
For

\[
f:\mathbb R^n\to\mathbb R,
\]

one reverse sweep gives

\[
\nabla f
\]

with derivatives with respect to all inputs at once.

A loss measures the error of a model's predictions. In machine learning it is usually a scalar, even when the model has many parameters.

For a vector-valued function

\[
\mathbf f:\mathbb R^n\to\mathbb R^m,
\]

we must first choose an output seed $\mathbf w\in\mathbb R^m$. Reverse mode then computes

\[
{
J_{\mathbf f}^T\mathbf w
}.
\]

The \textbf{\textit{vector-Jacobian product (VJP)}} is $\mathbf w^TJ_{\mathbf f}$. With column-vector notation, we write its transpose $J_{\mathbf f}^T\mathbf w$. The output seed $\mathbf w$ assigns a starting adjoint to each output; equivalently, we differentiate the scalar $L=\mathbf w^T\mathbf f$.

To reconstruct the complete Jacobian, use one reverse sweep per output basis vector.

\section*{Reverse-mode exercises}
Use exactly the same three functions as in the forward-mode section.

\subsection*{Exercise 2.1 - Scalar input, scalar output}
For

\[
f(x)=\log(xe^x+3)
\]

at $x=1$:

\begin{enumerate}
\item
Write the function as a sequence of elementary intermediate variables.

\item
Perform the forward pass and record all intermediate values.

\item
Initialize all adjoints to zero, then set the output adjoint equal to $1$.

\item
Propagate all adjoints backwards.

\item
Determine $\bar x=f'(1)$.

\end{enumerate}
\subsection*{Exercise 2.2 - Three inputs, scalar output}
For

\[
g(x,y,z)=xy+\sin z+y^2
\]

at

\[
(x,y,z)=(1,2,0),
\]

use the intermediate variables

\[
a=xy,
\qquad
b=\sin z,
\qquad
c=y^2,
\]

\[
d=a+b+c
\]

The scalar output is $g=d$.

\begin{enumerate}
\item
Perform the forward pass and record $a,b,c,d$.

\item
Initialize all adjoints to zero, then set $\bar d=1$ and propagate backwards.

\item
Determine the gradient $\nabla g(1,2,0)$ and compare it with Exercise 1.2.

\item
Explain why the adjoint of $y$ receives two contributions.

\end{enumerate}
\subsection*{Exercise 2.3 - Three inputs, two outputs}
For

\[
\mathbf h(x,y,z) = \begin{pmatrix} xy+z\\ x^2+\sin y-z^2\end{pmatrix}
\]

at

\[
(x,y,z)=(1,0,1),
\]

reverse mode needs an output seed because the output is not scalar.

\begin{enumerate}
\item
First output. Use

\[
\mathbf w_1 = \begin{pmatrix}1\\0\end{pmatrix}.
\]

That is, propagate backwards only from $h_1$.

Compute

\[
J_{\mathbf h}^T\mathbf w_1.
\]

\item
Second output. Use

\[
\mathbf w_2 = \begin{pmatrix}0\\1\end{pmatrix}
\]

Compute

\[
J_{\mathbf h}^T\mathbf w_2
\]

\item
Full Jacobian. Use the two results to reconstruct the complete Jacobian.

\item
General output seed. Without calculating the Jacobian from scratch, use

\[
\mathbf w = \begin{pmatrix}3\\-1\end{pmatrix}
\]

to calculate

\[
J_{\mathbf h}^T\mathbf w
\]

\end{enumerate}
\clearpage
\section*{Part III - Short conceptual questions}
\subsection*{Question 1}
Why does setting $\varepsilon^2=0$ cause the coefficient of $\varepsilon$ to behave like a derivative?

\subsection*{Question 2}
For a function

\[
f:\mathbb R^{100}\to\mathbb R,
\]

how many basis-direction forward sweeps would be needed to obtain the full gradient? How many reverse sweeps?

\subsection*{Question 3}
For a function

\[
f:\mathbb R\to\mathbb R^{100},
\]

which mode would naturally be more efficient for constructing the full Jacobian?

\subsection*{Question 4}
What is an adjoint $\bar v$?

\subsection*{Question 5}
Why do reverse-mode updates use accumulation such as

\[
\bar x\mathrel{+}=\cdots
\]

instead of simply assigning one value to $\bar x$?

\subsubsection*{Key concepts}
\begingroup
\small
\setlength{\tabcolsep}{4pt}
\begin{longtable}{p{\dimexpr 0.30\textwidth-2\tabcolsep\relax}p{\dimexpr 0.60\textwidth-2\tabcolsep\relax}}
\hline
Concept & Meaning in this exercise \\
\hline
\endfirsthead
Concept & Meaning in this exercise \\
\hline
\endhead
\textbf{Automatic differentiation (AD)} & Computing derivatives by applying the chain rule to elementary operations. \\[1.5pt]
\textbf{Loss} & A scalar measure of how poorly a model fits its target. \\[1.5pt]
\textbf{Dual number} & $x+\dot x\varepsilon$, where $\varepsilon\ne0$ and $\varepsilon^2=0$. \\[1.5pt]
\textbf{Input seed} & The initial tangent or input direction chosen for a forward sweep. \\[1.5pt]
\textbf{Directional derivative} & Rate of change along $\mathbf x+t\mathbf v$; for a scalar output, $\nabla g^T\mathbf v$. \\[1.5pt]
\textbf{Jacobian matrix} & Matrix of first partial derivatives, with outputs as rows and inputs as columns. \\[1.5pt]
\textbf{Forward mode / JVP} & Propagates values and tangents to compute $J\mathbf v$. \\[1.5pt]
\textbf{Sweep / pass} & One traversal of the calculation in a given direction. \\[1.5pt]
\textbf{Computational graph} & Nodes and directed dependencies representing a calculation. \\[1.5pt]
\textbf{Adjoint} & $\bar v=\partial L/\partial v$: sensitivity of the chosen scalar output $L$ to $v$. \\[1.5pt]
\textbf{Output seed} & Initial output adjoints; $1$ for a scalar output, or weights $\mathbf w$ for several outputs. \\[1.5pt]
\textbf{Reverse mode / VJP} & Propagates adjoints backwards to compute $J^T\mathbf w$, the transpose of $\mathbf w^TJ$. \\[1.5pt]
\textbf{Accumulation} & Adding all contributions to an adjoint when a variable affects the output through several paths. \\[1.5pt]
\hline
\end{longtable}
\endgroup

\end{document}
